% !TeX root = Thesis.tex

\chapter{Paper Overview and Problem Statement}\label{ch:introduction}

% \section{Soft Robotics}\label{sec:soft-robotics}

% todo read https://www.researchgate.net/publication/311479970_Soft_robotics_Technologies_and_systems_pushing_the_boundaries_of_robot_abilities once fulltext is granted

% \section{Overview of the paper}\label{sec:overview-of-the-paper}

This report gives an overview of the paper ``An Origami Continuum Robot Capable of Precise Motion Through Torsionally Stiff Body and Smooth Inverse Kinematics'' by Junius Santoso and Cagdas D. Onal from the Worcester Polytechnic Institute, Massachusetts\cite{santoso_origami_2021}.

\begin{wrapfigure}{r}{0.35\textwidth}
    \centering
    \includegraphics[width=0.3\textwidth]{images/robot}
    \caption{Continuum robotic arm consiting of $3$ modules\cite{santoso_origami_2021}.}
    \label{fig:arm}
\end{wrapfigure}

The authors identified precision to be the main challenge for continuum robot arms in related work\cite{santoso_origami_2021}.
They claim that missing torsional twist-resistance is one of the key causes for this precision loss.
Therefore, in this paper\cite{santoso_origami_2021}, the authors present a continuum robot arm with inherent passive torsional twist-resistance.
To achieve this, the robot consists of multiple segments of an origami construction.
An illustration of a continuum robotic arm can be seen in\ \prettyref{fig:arm}.
Origami constructions come with the advantage that they can be created by hand from sheets of readily-available materials like paper or plastic, which makes their production low-cost.
At the same time, they are lightweight and compliant to a degree that can be adjusted by the folding technique used.
This approach for designing soft robots is not new, an overview of previous works can be found in\cite{rus_design_2018}, as the authors point out\cite{santoso_origami_2021}.
Each segment of the robot is actuated by motors, which pull on cables that run through the origami-fold and are fixed to corresponding points on opposing end-plates.
By pulling on these cables, each segment can curve in any direction.
Multiple segments can be chained together and achieve complex motions in 3d-space.
The authors proposed minimally invasive surgery\cite{cianchetti_soft_2014}, rehabilitation\cite{polygerinos_soft_2015} and search and rescue\cite{tolley_resilient_2014} as areas of application for the continuum robot arm\cite{santoso_origami_2021}.
Although it is clear that the robot presented in the paper is not ready for deployment in any of these areas, the paper proposes improvements over previously presented continuum robot arms that take the state-of-the-art closer to such a real-world deployment.

This report will first go into more detail on the construction of the robot arm.
After that, the kinematics used to control the robot will be presented, as well as the implemented feedback-control.
Throughout these sections, corresponding experiments and results conducted by the authors will be summarized.
Finally, various areas of application, as well as the strengths and weaknesses of the presented design, will be discussed.
